\input{../tex-inputs/latex-leseansicht-vorspann}

\section[Arthur Schnitzler an Gustav Schwarzkopf, 16. 9. 1925]{L04410 Arthur Schnitzler an Gustav Schwarzkopf, 16. 9. 1925}
\nopagebreak\mylabel{L04410v}
\rehead{ }\normalsize\beginnumbering\briefempfaengerindex{Schwarzkopf, Gustav@\textsc{Schwarzkopf, Gustav}!zzzSchnitzler, Arthur@\emph{von Arthur Schnitzler}!1925-09-152@{16. 9. 1925}|(be}
\toendnotes[C]{\smallbreak\pagebreak[2]}
\correspDesc{Versand  durch Arthur Schnitzler am 16. 9. 1925 in Florenz
\newline{}Übermittlung  am 17. 9. 1925 in Florenz
\newline{}Erhalt  durch Gustav Schwarzkopf im Zeitraum [16. 9. 1925
                  – 20. 9. 1925?] in Wien}\toendnotes[C]{\smallbreak}
\Standort{DLA, A:Schnitzler, HS.1985.1.1897.}
\physDesc{Bildpostkarte, 85 Zeichen
\newline{}Handschrift: Bleistift, lateinische Kurrent
\newline{}Versand: Stempel: »\nobreak{}\oindex{Bahnhof Firenze Santa Maria Novella@\textbf{Bahnhof Firenze Santa Maria Novella}|pwk}Firenze Ferrovia, 17–IX 1925, 1–2\nobreak{}«.  }\toendnotes[C]{\smallbreak}\pstart{}{\pb}Austria\oindex{Österreich@\textbf{Österreich}|pw}\pend{}\pstart{}Hr Guſtav Schwarzkopf\pend{}\pstart{}Wien I\oindex{I., Innere Stadt@\textbf{I., Innere Stadt}|pw}\pend{}\pstart{}Tiefer Graben 17\oindex{Wien@\textbf{Wien}!I., Innere Stadt@\textbf{I., Innere Stadt}!Tiefer Graben 17@\textbf{Tiefer Graben 17}|pw}.\pend{}{\bigskip}
\pstart
           {\pb}\textcolor{gray}{\textbf{Firenze\oindex{Florenz@\textbf{Florenz}|pw}}}\hfill \textcolor{gray}{\textbf{Palazzo Pretorio o del Podestà\oindex{Museo Nazionale del Bargello@\textbf{Museo Nazionale del Bargello}|pw} ora Museo Nazionale\oindex{Museo Nazionale del Bargello@\textbf{Museo Nazionale del Bargello}|pwv}}}\pend
           \vspace{1em}
\pstart
           {\pb}16. 9. 25\pend
           \vspace{0.5em}
\pstart
           Herzliche Grüße!\pend
           
\pstart
           Ihr{\\[\baselineskip]}\spacefill\mbox{Arthur}\pend
           \leftskip=0em{}\selectlanguage{ngerman}\endnumbering\briefempfaengerindex{Schwarzkopf, Gustav@\textsc{Schwarzkopf, Gustav}!zzzSchnitzler, Arthur@\emph{von Arthur Schnitzler}!1925-09-152@{16. 9. 1925}|)be}\mylabel{L04410h}
\begin{anhang}
\end{anhang}\newcommand{\dateiname}{L04410}\newcommand{\titel}{Arthur Schnitzler an Gustav Schwarzkopf, 16. 9. 1925}\newcommand{\editorInnen}{Herausgegeben von Jahnke, SelmaMüller, Martin Anton}\input{../tex-inputs/latex-leseansicht-abspann}
